% GENERATED FILE. Edit canonical lessons, then run npm run generate:books.
% canonical-source-hash: fnv1a64:1c9b8edc8b30ace0
% canonical-lessons: HI-C220-jhanda, HI-C220-bacpan, HI-C220-javani, HI-C220-burhapa, HI-C220-jivan

\chapter{A flag, Childhood, Youth, Old age, Life}
\label{ch:hi-a-flag-childhood-youth-old-age-life}
\hlchaptermodality{220}

\begin{chapteropening}
\textbf{By the end of this chapter:} \emph{I can say a flag, childhood, youth, old age and life.}

Complete the last lesson of chapter 220: i can say a flag, childhood, youth, old age and life.
\end{chapteropening}

\section[\texorpdfstring{\hi{झंडा} (jhaṇḍā)}{झंडा (jhaṇḍā)}]{\hi{झंडा} (jhaṇḍā) — a flag}

\label{lesson:HI-C220-jhanda}



\begin{quote}
Before the new one: say the Hindi for a war, then the Hindi for an army.
\end{quote}

\subsection*{You'll want to know: \hi{झंडा}}
\textbf{\hi{झंडा}} — \emph{jhaṇḍā} — \textquotedblleft{}a flag\textquotedblright{}.

\begin{grammarlens}[title={What you've built}]
One more word for this chapter.
\end{grammarlens}

\subsection*{Your turn}
\begin{itemize}
\raggedright
  \item \emph{Say it:} \emph{jhaṇḍā}
  \item \emph{Say it:} \emph{jhaṇḍā}, once more
  \item \emph{Say it:} say \emph{jhaṇḍā}
  \item \emph{Recall:} say the Hindi for a war, then the Hindi for an army, then say \emph{jhaṇḍā} again
\end{itemize}

\begin{tcolorbox}[breakable,colback=teal!4,colframe=teal!35!black,title={Before you move on}]
What does \hi{झंडा} mean? (\textquotedblleft{}A flag\textquotedblright{}.) Say it once more.
\end{tcolorbox}
\section[\texorpdfstring{\hi{बचपन} (bacpan)}{बचपन (bacpan)}]{\hi{बचपन} (bacpan) — childhood}

\label{lesson:HI-C220-bacpan}



\begin{quote}
Before the new one: say the Hindi for an army, then the Hindi for a flag.
\end{quote}

\subsection*{You'll want to know: \hi{बचपन}}
\textbf{\hi{बचपन}} — \emph{bacpan} — \textquotedblleft{}childhood\textquotedblright{}.

\begin{grammarlens}[title={What you've built}]
One more word for this chapter.
\end{grammarlens}

\subsection*{Your turn}
\begin{itemize}
\raggedright
  \item \emph{Say it:} \emph{bacpan}
  \item \emph{Say it:} \emph{bacpan}, once more
  \item \emph{Say it:} say \emph{bacpan}
  \item \emph{Recall:} say the Hindi for an army, then the Hindi for a flag, then say \emph{bacpan} again
\end{itemize}

\begin{tcolorbox}[breakable,colback=teal!4,colframe=teal!35!black,title={Before you move on}]
What does \hi{बचपन} mean? (\textquotedblleft{}Childhood\textquotedblright{}.) Say it once more.
\end{tcolorbox}
\section[\texorpdfstring{\hi{जवानी} (javānī)}{जवानी (javānī)}]{\hi{जवानी} (javānī) — youth}

\label{lesson:HI-C220-javani}



\begin{quote}
Before the new one: say the Hindi for a flag, then the Hindi for childhood.
\end{quote}

\subsection*{You'll want to know: \hi{जवानी}}
\textbf{\hi{जवानी}} — \emph{javānī} — \textquotedblleft{}youth\textquotedblright{}.

\begin{grammarlens}[title={What you've built}]
One more word for this chapter.
\end{grammarlens}

\subsection*{Your turn}
\begin{itemize}
\raggedright
  \item \emph{Say it:} \emph{javānī}
  \item \emph{Say it:} \emph{javānī}, once more
  \item \emph{Say it:} say \emph{javānī}
  \item \emph{Recall:} say the Hindi for a flag, then the Hindi for childhood, then say \emph{javānī} again
\end{itemize}

\begin{tcolorbox}[breakable,colback=teal!4,colframe=teal!35!black,title={Before you move on}]
What does \hi{जवानी} mean? (\textquotedblleft{}Youth\textquotedblright{}.) Say it once more.
\end{tcolorbox}
\section[\texorpdfstring{\hi{बुढ़ापा} (buṛhāpā)}{बुढ़ापा (buṛhāpā)}]{\hi{बुढ़ापा} (buṛhāpā) — old age}

\label{lesson:HI-C220-burhapa}



\begin{quote}
Before the new one: say the Hindi for childhood, then the Hindi for youth.
\end{quote}

\subsection*{You'll want to know: \hi{बुढ़ापा}}
\textbf{\hi{बुढ़ापा}} — \emph{buṛhāpā} — \textquotedblleft{}old age\textquotedblright{}.

\begin{grammarlens}[title={What you've built}]
One more word for this chapter.
\end{grammarlens}

\subsection*{Your turn}
\begin{itemize}
\raggedright
  \item \emph{Say it:} \emph{buṛhāpā}
  \item \emph{Say it:} \emph{buṛhāpā}, once more
  \item \emph{Say it:} say \emph{buṛhāpā}
  \item \emph{Recall:} say the Hindi for childhood, then the Hindi for youth, then say \emph{buṛhāpā} again
\end{itemize}

\begin{tcolorbox}[breakable,colback=teal!4,colframe=teal!35!black,title={Before you move on}]
What does \hi{बुढ़ापा} mean? (\textquotedblleft{}Old age\textquotedblright{}.) Say it once more.
\end{tcolorbox}
\section[\texorpdfstring{\hi{जीवन} (jīvan)}{जीवन (jīvan)}]{\hi{जीवन} (jīvan) — life}

\label{lesson:HI-C220-jivan}



\begin{quote}
Before the new one: say the Hindi for youth, then the Hindi for old age.
\end{quote}

\subsection*{You'll want to know: \hi{जीवन}}
\textbf{\hi{जीवन}} — \emph{jīvan} — \textquotedblleft{}life\textquotedblright{}.

\begin{grammarlens}[title={What you've built}]
One more word for this chapter.
\end{grammarlens}

\subsection*{Your turn}
\begin{itemize}
\raggedright
  \item \emph{Say it:} \emph{jīvan}
  \item \emph{Say it:} \emph{jīvan}, once more
  \item \emph{Say it:} say \emph{jīvan}
  \item \emph{Recall:} say the Hindi for youth, then the Hindi for old age, then say \emph{jīvan} again
\end{itemize}

\begin{tcolorbox}[breakable,colback=teal!4,colframe=teal!35!black,title={Before you move on}]
What does \hi{जीवन} mean? (\textquotedblleft{}Life\textquotedblright{}.) Say it once more.
\end{tcolorbox}
